\section{Method}
\label{sec:method}

\paragraph{Notation.} $N$ is a parameter count (defined by a convention below), $D$ the number of training tokens, $C$ the training compute in FLOPs, and $L$ the validation cross-entropy loss in nats per token. A starred quantity ($N^*$, $D^*$) is compute-optimal for a given $C$.

\subsection{Models and compute accounting}
We train decoder-only Transformers in the GPT-2 style~\citep{radford2019gpt2}: pre-LayerNorm blocks, GELU MLP with 4$\times$ expansion, learned positional embeddings, no biases in linear layers, untied input and output embeddings, context length 512, and a GPT-2 byte-pair-encoding vocabulary padded to 50{,}304 tokens. We use seven model sizes (Table~\ref{tab:sizes}) with head dimension 64.

Training FLOPs per token are computed analytically as
\begin{equation}
  f(N) = 6\left(12\,n_\text{layer}\,d_\text{model}^2 + d_\text{model} V\right) + 6\,n_\text{layer}\,T\,d_\text{model},
\end{equation}
where $V$ is the vocabulary size and $T$ the context length; the first term counts matrix multiplications in the blocks and the output head (forward plus backward), the second counts causal attention. Embedding lookups cost no FLOPs. The total compute of a run is $C = D\,f$. Note that $f$ is \emph{not} proportional to any single parameter count here, because the output head contributes FLOPs but not block parameters.

\subsection{Parameter-count conventions}
\label{sec:conventions}
We compare three conventions: (\textbf{total}) all parameters; (\textbf{non-embedding}) transformer blocks and final LayerNorm only; (\textbf{non-embedding + head}) everything except the input token and position embeddings. These conventions are ours. They differ from the three interpretations of Schaeffer et al.~\citep{schaeffer2025}, which are three ways of computing the sizes of Chinchilla's models (the published counts, a formula assuming tied embeddings and no positional embeddings, and a best-fit variant of the attention term) that differ by at most 15.2\%. We compare embedding-based conventions because embedding parameters are the dominant ambiguity at our scale, and because including or excluding them acts like an additive offset to $N$, one of the perturbations Schaeffer et al.\ study.

\subsection{Training}
Each run uses AdamW~\citep{loshchilov2019adamw} with $\beta=(0.9, 0.95)$, weight decay 0.1 on matrices, gradient clipping at 1.0, bfloat16 autocast, and 32{,}768 tokens per optimizer step (gradient accumulation over micro-batches of 16 sequences). The learning rate follows linear warm-up (between 10 and 200 steps) and cosine decay to 10\% of the peak, scheduled over each run's own length, which is the setting IsoFLOP-style fits assume. The peak learning rate is set by the untuned heuristic $\min(2\times10^{-3},\,6\times10^{-4}\cdot 768/d_\text{model})$. Training data are consumed once in order, so no token is repeated. Each configuration is trained with a single random seed.

\subsection{Estimating compute-optimal allocations}
\label{sec:fits}
For each convention we apply the three approaches of Hoffmann et al.\ to the same runs. Every approach yields a power law $N^* \propto C^{a}$ and $D^* \propto C^{b}$; we additionally report $N^*$ and $D^*$ at a reference budget $C_\text{ref}$, the geometric mean of the budgets in the grid.

\paragraph{Approach 1: training-curve envelope.} Each run's validation loss is recorded about eight times during training. We interpolate each run's curve linearly in $\log C$ and, at every $C$ on a dense grid where at least two model sizes are available, select the run with the lowest loss. We regress $\log N$ and $\log D$ of the selected runs on $\log C$. Because intermediate points come from runs whose learning rate has not finished decaying, and because the choice is among discrete model sizes, this approach is coarse.

\paragraph{Approach 2: IsoFLOP profiles.} For every budget with at least three runs, we fit a parabola to loss as a function of $\log N$ and take its minimum as $N^*$ (requiring an interior minimum); $D^*$ is interpolated from the runs at the same budget. We then regress $\log N^*$ and $\log D^*$ on $\log C$.

\paragraph{Approach 3: parametric loss surface.} We fit $L(N,D) = E + A N^{-\alpha} + B D^{-\beta}$ to all runs by minimizing the Huber loss ($\delta = 10^{-3}$) of the log-residuals with L-BFGS-B from 162 initializations. For a budget $C$, the optimum minimizes $L(N, C/f(N))$ over $N$, where $f(N)$ is a log-log linear fit of FLOPs per token against $N$ across the model ladder. This makes the optimum consistent with the actual compute accounting for every convention.

\subsection{Robustness metrics and uncertainty}
\label{sec:robust}
For each approach we report, across conventions, the spread of the exponent $a$ (maximum minus minimum) and the relative spread of $D^*(C_\text{ref})$, which is a physical quantity (a token count) and therefore comparable across conventions, whereas $N^*$ and $a$ depend on how $N$ is defined. Pairwise gaps between approaches address RQ3. Approach~1 selects runs by loss alone, so its $D^*$ does not depend on the convention by construction.

We quantify fitting uncertainty with a bootstrap (\TODO{200} resamples): runs are resampled with replacement within each budget (Approaches 1 and 2) or overall (Approach 3), and we report percentile 95\% intervals. These intervals capture variation due to which runs are included, \emph{not} variation due to training seeds.

\subsection{Structured perturbations of the parameter counts}
\label{sec:perturb}
Following Schaeffer et al.~\citep{schaeffer2025}, we perturb the parameter counts of the total-parameter convention in four ways and refit all three approaches each time, leaving $D$ and the losses untouched: (i) a multiplicative constant, $N\to cN$, $c\in\{0.1,0.3,1,3,10\}$; (ii) an additive constant, $N\to N+c$, $c\in\{-10,0,10,40\}$M; (iii) a systematic bias that stretches $\log N$ around its geometric mean with slope $s\in\{0.5,0.75,1,1.5,2\}$; and (iv) log-normal noise, $N\to e^{\delta}N$ with $\delta\sim\mathcal N(0,\sigma^2)$, $\sigma\in\{0.01,0.1,0.3\}$, averaged over 10 draws. Because a counting error is a property of the architecture, each perturbation is applied per model size, so all runs of one size share it. Where Schaeffer et al.\ assume $C\approx cND$ with an unspecified constant, we keep the exact FLOPs accounting $f(N)$, recomputed from the perturbed counts. We report the optimal tokens-per-parameter ratio and the exponent $a$ at each level.

\paragraph{Grid-shrinking sweep (RQ2).} We repeat the analysis after removing each budget in turn, each model size in turn, the two extreme sizes, and every other size, and track the spread across conventions as a function of the number of runs.

\paragraph{Validation on a known law.} To check the fitting code, we generate synthetic runs from a known law and verify that all three approaches recover its optimal allocation (see Appendix).
